


%%%%%%%%%%%%%%%%%%%%%%%%%%%%%%%%%%%%%%%%%%%%%%%%%%%%%%%%%%%%%%%%%%%%%%%%%%%%%%%%
%%%%%%%%%%%%%%%%%%%%%%%%%%%%%%%%%%%%%%%%%%%%%%%%%%%%%%%%%%%%%%%%%%%%%%%%%%%%%%%%


\section{Preliminaries}\label{sec:zinc-preliminaries}
\label{s:extended_preliminaries}


% \albert{Should unify usage of boldface and non-boldface letters for vectors and matrices across the paper.}

\subsection{Basic notions}\label{s:basic_notions}




For $k\geq 1$, let $[k]\defeq \{1,\ldots, k\}$, and let $\range{0}{k}$ be $\{0\}\cup[k]$. We write $x\getsr S$ for uniform sampling from a finite set $S$, and $x\gets\mathcal{D}$ for sampling from a distribution $\mathcal{D}$.


%\albert{we are not doing that?} Vectors and tuples are denoted by boldface letters $\vv, \uu, \ldots$, with $\vv_i$ the $i$-th entry. When $\vv$ denotes a matrix, $\vv_{i}$ is the $i$-th row and $\vv_{i,j}$ the $(i,j)$-th entry.
%For execution traces (cf.\ \cref{s:uair}), we write $\vv_{\colindex}[\rowindex]$ for the entry at row $\rowindex$ and column $\colindex$, and $\vv[\rowindex]$ for the row vector at row $\rowindex$.
%Given a vector $\vv$ of length $n$, we write $\set(\vv)\defeq\{\vv_i:i\in[n]\}$ for the set of entries of~$\vv$. %For an $m\times n$ matrix $M$ and a column index $j\in[n]$, we write $\set(M_{\cdot,j})\defeq\{M_{i,j}:i\in[m]\}$ for the set of entries in column~$j$.



%%%%%%%%%%%%%%%%%%%%%%%%%%%%%%%%%%%%%%%%%%%%%%%%%%%%%%%%%%%%%%%%%%%%%%%%%%%%%%%%
%%%%%%%%%%%%%%%%%%%%%%%%%%%%%%%%%%%%%%%%%%%%%%%%%%%%%%%%%%%%%%%%%%%%%%%%%%%%%%%%

We use $\KK$ to denote (possibly infinite) fields and $R$ to denote rings. We write $\FF_q$ for the finite field with $q$ elements (where $q$ is a prime or a prime power), $\QQ$ for the field of rational numbers, and $\ZZ$ for the ring of integers. We use standard terminology and notation: an ideal $I$ generated by a set $S\subseteq R$ is denoted $(S)$, the quotient ring is denoted $R/I$.   Sometimes we write $(S)_R$ to emphasize that the ideal is an ideal of the ring $R$. A ring is an integral domain if the product of any two nonzero elements is nonzero.



\parhead{Polynomial rings and commutative algebra}
For variables $\YY=(Y_1,\ldots, Y_\mu)$, $R[\YY]$ denotes multivariate polynomials over $R$, and $R^{\multilin}[\YY]$ the multilinear polynomials (individual degree $\le 1$ in each variable).

We often work with multivariate polynomials whose coefficients are themselves univariate polynomials. For a ring $\cR$ (typically $\ZZ[X]$, $\QQ[X]$, or $\local{q}[X]$) and variables $\YY=(Y_1,\ldots, Y_\mu)$, we write $\cR[\YY]$ for polynomials in $\YY$ with coefficients in $\cR$. We view elements of $(\ZZ[X])[\YY]$ as polynomials in $\YY$ (not as elements of $\ZZ[X,\YY]$), denoting them by $f(\YY)$.
For polynomial rings, $\cR^{<d, <B}[X]$ denotes polynomials of degree $<d$ with coefficients in $\cR^{<B}$; we may write $\cR^{<d}[X]$ when the bit bound is clear.
Ring homomorphisms $f:\cR_1 \to \cR_2$ extend coefficient-wise to $\cR_1[\YY] \to \cR_2[\YY]$.
More generally, all maps are extended entrywise to vectors and matrices.


For any field $\FF_{p^e}$ containing subfield $L$ we can map elements of $L$ into $\FF_{p^e}$ via the canonical embedding $\embedding_{p^{e}}: L \hookrightarrow \FF_{p^e}$. We extend this embedding to the ring $\FF_{p^e}[X]$ its subrings $L[X]$ by applying $\embedding_{p^{e}}$ coefficient-wise. If $(\generator) = \Ideal \subseteq \FF_q[X]$ is an ideal, we note $(\embedding_{q^\ell}(\generator)) = \Ideal' \subseteq \FF_{q^\ell}[X]$ is also.% \psiv{Todo: confirm this and that it's the relevant statement for the IOR.}



\label{sec:zinc-local-rings}
The subset $\loc{p} = \{a/b\in \QQ : p\nmid b\}$ of $\QQ$ is a ring called the \emph{localization of $\ZZ$ at the prime~$p$}. This ring admits a surjective homomorphism $\phi_p: \loc{p} \to \FF_p$, $a/b \mapsto a \cdot b^{-1}$, with kernel $\ker(\phi_p) = p\loc{p}$.


For prime powers $q=p_\primeindex^{e_\primeindex}$, we set $\loc{q_\primeindex}\defeq \loc{p_\primeindex}$, and extend $\phi$ by defining $\phi_{q_\primeindex} = \iota_{q_\primeindex} \circ \phi_{p_\primeindex}$. The \emph{intersection localization ring} is
\begin{equation}\label{e: multi_localization}
\loc{q_1\cdots q_{n}}
\;\defeq\;\bigcap_{\primeindex\in[n]}\loc{p_\primeindex}
\;=\;
\Bigl\{a/b\in \QQ \ \bigm|\ \gcd\!\bigl(b,\textstyle\prod_{\primeindex} p_\primeindex\bigr)=1\Bigr\}.
\end{equation}


\parhead{Bitstring representation}
\label{sec:bitstring-rep}
We use standard representations: integers are represented in two's complement, so an integer $n\in\ZZ$ has bit-size $\lceil\log_2(|n|+1)\rceil + 1$; rationals $a/b\in\QQ$ are represented as pairs $(a,b)$ in \emph{lowest terms} (i.e., $\gcd(a,b)=1$ and $b>0$), with bit-size the sum of the bit-sizes of $a$ and $b$; polynomials are represented as tuples of coefficients (in increasing degree order), with bit-size the sum of the coefficient bit-sizes.
We write $\ZZ^{<B}$, $\QQ^{<B}$, $\ZZ^{<d,<B}[X]$, $\QQ^{<d,<B}[X]$, and $\local{q}^{<d,<B}[X]$ for elements with bit-size less than $B$ (and degree $<d$ for polynomials).
For a ring $\cR$, we write $\cR^{<B}$ for elements representable with fewer than $B$ bits.


\parhead{Multilinear extensions (MLE)}
\label{sec:multilinear-extensions}
A multilinear polynomial $f(\YY)\in R^{\multilin}[\YY]$ is uniquely determined by its values on $\BB^\mu$. For any map $f: \BB^\mu \to R$, we write $\multilinmle{f}(\YY)$ for its \emph{multilinear extension} (MLE), the unique multilinear polynomial that agrees with $f$ on the Boolean hypercube. 
For a vector $\mathbf{v}\in R^{2^\mu}$, we write $\multilinmle{\mathbf{v}}(\YY)$ for the MLE of the map $\yy\mapsto v_\yy$ (interpreting $\yy\in\BB^\mu$ as an integer in binary).
We also use the standard \emph{equality polynomial}
\begin{equation}\label{e: MLE}
\meq(\yy; \YY) \defeq \prod_{i\in [\mu]} \left(y_i Y_i + (1-y_i)(1-Y_i)\right),
\end{equation}
which satisfies $\meq(\yy;\zz)=1$ if $\yy=\zz$ and $\meq(\yy;\zz)=0$ otherwise for $\yy,\zz\in\BB^\mu$. 


\parhead{Schwartz--Zippel lemma for exceptional sets}
\label{sec:algebraic-tools}
We will need the following general form of Schwartz-Zippel lemma in order to address certain ideal-membership constraints in our proof system. A subset $S\subseteq \cR$ is \emph{exceptional} if $a-b$ has a multiplicative inverse for all distinct $a,b\in S$.
\begin{lemma}[Schwartz--Zippel lemma for exceptional sets]
\label{lem:sz_exceptional_sets} \albert{Reference somewhere in the literature for a proof }
Let $S\subseteq\cR$ be a finite exceptional set and $f(\YY)\in \cR[\YY]$ a nonzero multilinear polynomial in $\mu$ variables. Then
$\Pr_{\rr\getsr S^\mu}[f(\rr)=0] \le \mu/|S|$.
\end{lemma}

% \subsection{Interactive Oracle Proofs (IOPs), and Polynomial IOPs (PIOPs)}\label{sec:IOPs}


\subsection{Indexed relations and oracles.}
\label{sec:zinc-polynomial-commitments}


An \emph{indexed relation} $\REL$ is a set of index-input-witness triples $(\idx,\inp;\wit)$. The \emph{language} associated with $\REL$ is
\(
  \LANG(\REL)\ \defeq\ \bigl\{(\idx,\inp)\ \bigm|\ \exists\,\wit \text{ such that } (\idx,\inp;\wit)\in\REL\bigr\}.
\)



An \emph{oracle} to a string or vector $\pi$ of elements from a ring $R$ is a black-box that allows querying any position $j$ and obtaining $\pi_j$. We denote it by $\oracle{\pi}$. 
An \emph{oracle to a polynomial} $f(\XX)$ over $R$ is a black-box that allows querying any point $\rr$ in the domain of $f$ and obtaining $f(\rr)$. We denote it by $\oracle{f}$.

% We provide standard definitions for IOPs, PIOPs, and their properties in \cref{sec:IOPs}. 


% \begin{definition}[Polynomial Interactive Oracle Proof (PIOP)]
% \label{def:piop}
% A \emph{polynomial interactive oracle proof (PIOP)} for an indexed relation $\REL$ is a public-coin interactive protocol $(\prover,\verifier)$ in which:
% \begin{itemize}
%   \item The index $\idx$ and input $\inp$ may contain oracles to polynomials over a ring $R$ with prescribed number of variables, degrees, and coefficient domains.
%   \item In each round, the prover $\prover$ may send, among other data, oracles to polynomials (with prescribed structure), and the verifier $\verifier$ sends a uniformly random challenge.
%   \item At the end of the interaction, the verifier queries the oracles (from the index, input, and prover messages) at points of its choosing and outputs $\mathsf{accept}$ or $\mathsf{reject}$.
% \end{itemize}
% \noindent\textbf{Completeness.} For every $(\idx,\inp;\wit)\in\REL$, 
% \[
%   \Pr\bigl[\langle\prover(\idx,\inp;\wit),\verifier(\idx,\inp)\rangle = \mathsf{accept}\bigr]\ \ge\ 1-\varepsilon_{\mathsf{cmp}},
% \]
% where $\varepsilon_{\mathsf{cmp}}$ is the \emph{completeness error}. When $\varepsilon_{\mathsf{cmp}}=0$, the protocol has \emph{perfect completeness}.

% \noindent\textbf{Soundness.} For every $(\idx,\inp)\notin\LANG(\REL)$ and every (computationally unbounded) prover $\prover^*$ that sends correctly typed oracle messages (cf.\ \cref{rem:zinc-oracle_typing}),
% \[
%   \Pr\bigl[\langle\prover^*,\verifier(\idx,\inp)\rangle = \mathsf{accept}\bigr]\ \le\ \varepsilon_{\mathsf{snd}},
% \]
% where $\varepsilon_{\mathsf{snd}}$ is the \emph{soundness error}.

% \noindent\textbf{Knowledge soundness.} \albert{Should be round-by-round knowledge soundness? Let's wait for Psi's section  } The PIOP is \emph{knowledge sound} with knowledge error $\knowledgeerror$ if there exists a polynomial-time \emph{extractor} $\extractor$ such that for every prover $\prover^*$ that sends correctly typed oracle messages (cf.\ \cref{rem:zinc-oracle_typing}), every index $\idx$, and every input $\inp$,
% \[
%   \Pr\Bigl[
%     \langle\prover^*,\verifier(\idx,\inp)\rangle = \mathsf{accept}
%     \ \wedge\
%     (\idx,\inp;\extractor^{\prover^*}(\idx,\inp))\notin\REL
%   \Bigr]
%   \ \le\ \knowledgeerror,
% \]
% where $\extractor^{\prover^*}$ denotes that the extractor has oracle access to the prover's messages.
% \end{definition}



\begin{remark}[Oracle typing assumption]
\label{rem:zinc-oracle_typing}
Soundness and knowledge soundness for PIOPs are required to hold only for malicious provers $\prover^*$ that send \emph{correctly typed} oracles---i.e., oracles that actually correspond to polynomials with the prescribed number of variables, degrees, and coefficient domains. Informally speaking, enforcing that a (computationally bounded) prover sends such oracles is the responsibility of the polynomial commitment scheme (PCS) or other cryptographic mechanisms (such as IOPP) used to compile the PIOP into a concrete argument system.
\end{remark}


% \begin{definition}[Interactive Oracle Proof (IOP)]
% \label{def:iop}
% An \emph{interactive oracle proof (IOP)} is defined analogously to a PIOP (\cref{def:piop}), except that oracles are to strings (or vectors) of ring elements rather than to polynomials. The verifier queries individual positions of these strings.
% \end{definition}


% An \emph{interactive proof system} for $\REL$ is an IOP in which neither the index, input, nor prover messages contain any oracles. Completeness, soundness, and knowledge soundness are as in \cref{def:piop}. An \emph{interactive argument} is defined similarly, except that soundness is required to hold only against computationally bounded (PPT) provers.

% %\begin{lemma}[Sound PIOPs are knowledge sound -- Adaptation of {\cite[Lemma 2.3]{hyperplonk}}]\label{lem: piopks} \albert{We don't use this do we?}
% %Let $\prot$ be a PIOP for a relation $\REL$ with soundness error $\varepsilon$. Assume that, for all $(\idx, \inp; \wit)\in\REL$, $\wit$ consists only of multilinear polynomials such that $\inp$ contains oracles to these polynomials. Then the PIOP has knowledge soundness error $\varepsilon$, and the extractor runs in time $O(|\wit|)$.
% %\end{lemma}
% %The proof is given in \cref{app:lem_piopks_proof}.



% \subsection{Constrained linear codes, IOPs of Proximity (IOPPs), and Mutual Correlated Agreement (MCA)}\label{def:prelims_iopp}

% \parhead{Constrained linear codes}

% A \emph{linear code} over a field  $\KK$ is a linear subspace $\mc{C}\subseteq \KK^n$.
% The \emph{block length} is $n$, and the \emph{dimension} $k=\dim(\mc{C})$ is the dimension of the subspace.
% A \emph{generator matrix} $M\in \KK^{n\times k}$ for $\mc{C}$ is a matrix whose columns form a basis for $\mc{C}$, so $\mc{C}=\{Mw : w\in \KK^k\}$.
% \albert{Define distance first and then minimum distance and so on} The \emph{relative distance} $\delta\in(0,1]$ of $\mc{C}$ is the minimum, over all distinct $\cc,\cc'\in\mc{C}$, of the fraction of coordinates in which $\cc$ and $\cc'$ differ.
% For a vector $\vv\in \KK^n$, the \emph{relative Hamming distance} from $\vv$ to the code is $\dist(\vv,\mc{C})=\min_{\cc\in\mc{C}}\dist(\vv,\cc)$, where $\dist(\vv,\cc)$ denotes the fraction of coordinates in which $\vv$ and $\cc$ differ.

% Given a linear code $\mc{C}\subseteq \KK^n$ with generator matrix $M\in \KK^{n\times k}$, the \emph{$J$-wise interleaving} of $\mc{C}$ is the code
% \[
%   \mc{C}^J \;:=\; \{WM^T : W\in \KK^{J\times k}\} \;\subseteq\; \KK^{J\times n}.
% \]
% A codeword in $\mc{C}^J$ can be viewed as a tuple $(v_1,\ldots,v_J)$ of $J$ codewords $v_i\in\mc{C}$.


% \subparagraph{IOPs of proximity.}
% An \emph{IOP of proximity (IOPP)} is an IOP variant designed to test whether an input oracle is \emph{close} to a member of some set (typically a code or a constrained code), rather than exactly a member.

% \begin{definition}[Proximity relation]
% \label{def:proximity_relation}
% The \emph{proximity relation} $\REL_{\mathsf{prox}}$ for (possibly constrained) codes $\cC$ over $\KK$ and proximity parameters $\beta$ is defined as follows:
% \[
%   \REL_{\mathsf{prox}}
%   \;:=\;
%   \bigl\{(\idx,\inp;\wit)\ \bigm|\ \idx = (\mathcal{C},\beta),\ \inp = \oracle{\vv},\ \wit = \cc \in \mathcal{C},\ \dist(\vv,\cc) \le \beta\bigr\}.
% \]
% \end{definition}

% \begin{definition}[Interactive Oracle Proof of Proximity (IOPP)]
% \label{def:iopp}
% Let $\idx=(\cC, \beta)$ be an index for $\REL_{\mathsf{prox}}$. 
% An \emph{interactive oracle proof of proximity (IOPP)} for $\idx$ is an IOP for the proximity relation $\REL_{\mathsf{prox}}$ with  index restricted to $\idx = (\mathcal{C},\beta)$ (\cref{def:proximity_relation}), and with completeness only guaranteed when $\vv$ is in the code $\cC$, rather than $\beta$-close to it.
% \begin{comment}
% Concretely, it is a public-coin interactive protocol $(\prover,\verifier)$ in which:
% \begin{itemize}
%   \item $\prover$ and $\verifier$ receive $(\idx, \inp)$ as inputs, where  $\inp = \oracle{\vv}$ for some $\vv \in \KK^n$. The honest prover additionally receives a witness $\cc \in \mathcal{C}$ such that $\dist(\vv,\cc) \le \beta$.
%   \item In each round, $\prover$ sends oracles to strings, and $\verifier$ sends a uniformly random challenge.
%   \item At the end of the interaction, $\verifier$ queries the input oracle $\oracle{\vv}$ and the prover's oracle messages at positions of its choosing and outputs $\mathsf{accept}$ or $\mathsf{reject}$.
% \end{itemize}
% \noindent \textbf{Completeness.} For every $\vv \in \mathcal{C}$, there exists a prover $\prover$ such that
% \[
%   \Pr\bigl[\langle\prover(\idx,\inp,\wit),\verifier(\idx,\inp)\rangle = \mathsf{accept}\bigr]\ \ge\ 1-\varepsilon_{\mathsf{cmp}}.
% \]
% \noindent \textbf{Soundness.} For every $\vv \in \KK^n$ with $\dist(\vv, \mathcal{C}) > \beta$ and every prover $\prover^*$,
% \[
%   \Pr\bigl[\langle\prover^*,\verifier(\idx,\inp)\rangle = \mathsf{accept}\bigr]\ \le\ \varepsilon_{\mathsf{snd}}.
% \]
% \albert{Missing RBR Knowledge Soundness}
% \end{comment}
% \end{definition}



% \parhead{Mutual Correlated Agreement}
% \input{IOPP/mutual_correlated_agreement.tex}


\subsection{Polynomial Interactive Oracle Reductions (PIORs)}
\label{sec:iors}

Polynomial Interactive Oracle Reductions (PIORs) generalize PIOPs. A \emph{public-coin polynomial IOR} from a relation $\Relation$ to a relation $\Relation'$ is a pair $\Pi = (\prover,\verifier)$ specified as follows.

The prover $\prover$ is an interactive algorithm that receives as input $(\idx,\inp,\wit)\in \Relation$.  The verifier $\verifier$ is an interactive polynomial oracle algorithm begins with inputs $(\idx,\inp)$, which may include both polynomial oracles and explicit values. They interact for $k$ rounds. In round $i\in[k]$, the prover sends a polynomial oracle $\provermsg_i$ and the verifier sends a uniformly random message $\verifiermsg_i$.  After $k$ rounds of interaction, the verifier makes evaluation queries to the polynomial oracles in its inputs and sent in prover messages $\provermsg_1,\dots,\provermsg_k$ as oracles. Then:
\begin{itemize}
  \item $\verifier$ either rejects or outputs a new index-instance $(\idx',\inp')$, and
  \item $\prover$ outputs a new witness $\wit'$ such that $(\idx',\inp',\wit') \in \Relation'$.
\end{itemize}

Without loss of generality, we may view $\verifier$'s output as a deterministic function of its input and the interaction transcript. Note that oracles in $\inp'$ are specified implicitly in terms of oracles given in $\inp$ and $\provermsg_1,\dots,\provermsg_k$.

\begin{definition}[Completeness]
\label{def:ior-completeness}
We say that $\Pi$ has \emph{completeness} with error $\completenesserror$ if for every $(\idx,\inp,\wit)\in \Relation$, letting the interaction proceed honestly and writing $(\idx',\inp') \gets \verifier\bigl(\idx,\inp,\provermsg_1,\verifiermsg_1,\dots,\provermsg_k,\verifiermsg_k\bigr)$ and $(\provermsg_k,\wit') \gets \prover(\idx,\inp,\wit,\verifiermsg_1,\dots,\verifiermsg_{k-1})$,
  we have
  \[
    \Pr_{\verifiermsg_1,\dots,\verifiermsg_k}\left[\, (\idx',\inp',\wit') \not\in \Relation' \right] \leq \epsilon(\idx,\inp).
  \]
\end{definition}
We omit a standard soundness definition as we will prove the stronger notion of round-by-round knowledge soundness.% \psiv{Which hopefully implies Fiat--Shamir security of the corresponding non-interactive reduction in the ROM}

We will use the definition of Round-by-Round knowledge soundness from \cite{BCFW25}.
\begin{definition}[Knowledge state function]
\label{def:knowledge-state-function}
Let $\Pi$ be an IOR from a relation $\Relation$ to a relation $\Relation'$.
A \emph{knowledge state function} for $\Pi$ is a (possibly inefficient) function $\kstate$, that takes as input a statement $(\idx,\inp,\impinp)$, an interaction transcript $\tr$, a knowledge witness $\wk$, and outputs a bit. It must satisfy:
\begin{itemize}[nosep]
  \item \textbf{Empty transcript.} If $\tr=\emptyset$, then $\kstate(\idx,\inp,\impinp,\tr, \wk)=1$ if and only if $(\idx,\inp,\impinp,\wk)\in \Relation$.
  \item \textbf{Prover moves.} If $\tr$ is a transcript where the prover is about to move and $\kstate(\idx,\inp,\impinp,\tr, \wk)=0$, then for every prover message $\provermsg$,
  \[
    \kstate(\idx,\inp,\impinp,\tr\Vert \provermsg, \wk)=0.
  \]
  \item \textbf{Full transcript.} If $\tr=(\provermsg_1,\verifiermsg_1,\dots,\provermsg_k,\verifiermsg_k)$ is a full transcript, then $\kstate(\idx,\inp,\impinp,\tr, \wk)=1$ if and only if $\verifier^{\impinp,\provermsg_1,\dots,\provermsg_k}\bigl(\idx,\inp,\verifiermsg_1,\dots,\verifiermsg_k\bigr)$ outputs $(\inp',\impinp')$ such that $(\idx,\inp',\impinp',\wk)\in \Relation'$.
\end{itemize}
\end{definition}

\begin{definition}[Round-by-round knowledge soundness]
\label{def:rbr-knowledge-new}
An IOR $\Pi$ from $\Relation$ to $\Relation'$ has \emph{round-by-round knowledge soundness} with errors $(\knowledgeerror_1,\dots,\knowledgeerror_k)$
%and extraction times $(\extractiontime_1,\dots,\extractiontime_k)$
if there exists a knowledge state function $\kstate$ for $\Pi$ and a deterministic extractor $\extractor$ such that for every statement $(\idx,\inp,\impinp)$, every round index $i\in[k]$, and every interaction transcript where the verifier is about to move it holds:
%, the extractor $\extractor$ runs in time at most $\extractiontime_i$ and
\[
  \Pr_{\verifiermsg_i}\left[
    \begin{array}{l}
    \exists \wk=0 \text{ s.t. }\\
    \wedge\ \kstate(\idx,\inp,\impinp,\tr, \extractor(\idx,\inp,\impinp,\tr\Vert\verifiermsg_i, \wk))=1 \\
    \wedge\ \kstate(\idx,\inp,\impinp,\tr\Vert\verifiermsg_i, \wk)=1 \\
    \end{array}
  \right]
  \ \le\ \knowledgeerror_i(\idx,\inp,\impinp).
\]
% The \emph{total extraction time} is defined as $\sum_{i\in[k]} \extractiontime_i$.
\end{definition}

At times we will also use a prior, stronger notion of round-by-round knowledge soundness, which is known to imply the round-by-round knowledge soundness according to \cref{def:rbr-knowledge-new} with the same errors \cite{BCFW25}.

\begin{definition}[Knowledge state function ]
\label{def:knowledge-state-function-strong}
Let $\Pi$ be an IOR from a relation $\Relation$ to a relation $\Relation'$.
A \emph{knowledge state function} for $\Pi$ is a (possibly inefficient) function $\kstate$, that takes as input a statement $(\idx,\inp)$ and an interaction transcript $\tr$, and outputs a bit. It must satisfy:
\begin{itemize}[nosep]
  \item \textbf{Empty transcript.} If $\tr=\emptyset$, then $\kstate(\idx,\inp,\tr)=0$.
  \item \textbf{Prover moves.} If $\tr$ is a transcript where the prover is about to move and $\kstate(\idx,\inp,\tr)=0$, then for every prover message $\provermsg$,
  \[
    \kstate(\idx,\inp,\tr\Vert \provermsg)=0.
  \]
  \item \textbf{Full transcript.} If $\tr=(\provermsg_1,\verifiermsg_1,\dots,\provermsg_k,\verifiermsg_k)$ is a full transcript, then $\kstate(\idx,\inp,\tr)=1$ if and only if $\verifier\bigl(\idx,\inp,\provermsg_1,\verifiermsg_1,\dots,\provermsg_k,\verifiermsg_k\bigr)$ outputs $(\idx',\inp')$ such that there exists $\wit$' with $(\idx',\inp',\wit')\in \Relation'$.
\end{itemize}
\end{definition}

\begin{definition}[Round-by-round knowledge soundness]
\label{def:rbr-knowledge}
An IOR $\Pi$ from $\Relation$ to $\Relation'$ has \emph{round-by-round knowledge soundness} with errors $(\knowledgeerror_1,\dots,\knowledgeerror_k)$
%and extraction times $(\extractiontime_1,\dots,\extractiontime_k)$
if there exists a knowledge state function $\kstate$ for $\Pi$ and a deterministic extractor $\extractor$ such that for every statement $(\idx,\inp)$, every round index $i\in[k]$, and every interaction transcript where the verifier is about to move it holds:
%, the extractor $\extractor$ runs in time at most $\extractiontime_i$ and
\[
  \Pr_{\verifiermsg_i}\left[
    \begin{array}{l}
    \kstate(\idx,\inp,\tr)=0 \\
    \wedge\ \kstate(\idx,\inp,\tr\Vert\verifiermsg_i)=1 \\
    \wedge\ (\idx,\inp,\extractor(\idx,\inp,\tr)) \not\in \Relation
    \end{array}
  \right]
  \ \le\ \knowledgeerror_i(\idx,\inp).
\]
% The \emph{total extraction time} is defined as $\sum_{i\in[k]} \extractiontime_i$.
\end{definition}

% \psiv{Todo: prove that (our variant of) round-by-round knowledge soundness implies straightline state-restoration knowledge soundness}

% \parhead{Efficiency measures}
% We track several efficiency measures:
% \begin{itemize}
%   \item rounds $k$;
%   \item proof length $\ell$, the total number of nonzero coefficients across $\provermsg_1,\dots,\provermsg_k$;
%   \item input-query complexity $q_y$, the number of evaluations queried from all oracles
%   \item proof-query complexity $q_\pi$, the number of evaluations queried from $\provermsg_1,\dots,\provermsg_k$;
%   \item verifier randomness $r := \sum_{i=1}^k r_i$ (where $\verifiermsg_i\in\{0,1\}^{r_i}$);
%   \item prover time $\provertime$, measured in algebraic operations;
%   \item verifier time $\verifiertime$, measured in algebraic operations.
% \end{itemize}

